\begin{equation}
    |0\rangle = \begin{pmatrix} 1 \\ 0 \end{pmatrix}, \quad |1\rangle = \begin{pmatrix} 0 \\ 1 \end{pmatrix}
\end{equation}
\begin{itemize}
    \item $|0\rangle, |1\rangle$ : Vektor basis komputasi kuantum (\textit{computational basis states}).
\end{itemize}

\begin{equation}
    \sigma_x = \begin{pmatrix} 0 & 1 \\ 1 & 0 \end{pmatrix}, \quad \sigma_y = \begin{pmatrix} 0 & -i \\ i & 0 \end{pmatrix}, \quad \sigma_z = \begin{pmatrix} 1 & 0 \\ 0 & -1 \end{pmatrix}
\end{equation}
\begin{itemize}
    \item $\sigma_x, \sigma_y, \sigma_z$ : Matriks Pauli $X, Y,$ dan $Z$.
\end{itemize}

\begin{equation}
    \hat{R}_y(\theta) = \exp\left(-i \frac{\theta}{2} \sigma_y\right)
\end{equation}
\begin{equation}
    \hat{R}_y(\theta) = \hat{I} \cos\left(\frac{\theta}{2}\right) - i \sigma_y \sin\left(\frac{\theta}{2}\right)
\end{equation}
\begin{equation}
    \hat{R}_y(\theta) = \begin{pmatrix} 1 & 0 \\ 0 & 1 \end{pmatrix} \cos\left(\frac{\theta}{2}\right) - i \begin{pmatrix} 0 & -i \\ i & 0 \end{pmatrix} \sin\left(\frac{\theta}{2}\right)
\end{equation}
\begin{equation}
    \hat{R}_y(\theta) = \begin{pmatrix} \cos\left(\frac{\theta}{2}\right) & -\sin\left(\frac{\theta}{2}\right) \\ \sin\left(\frac{\theta}{2}\right) & \cos\left(\frac{\theta}{2}\right) \end{pmatrix}
\end{equation}

\begin{equation}
    \hat{R}_z(\theta) = \exp\left(-i \frac{\theta}{2} \sigma_z\right)
\end{equation}
\begin{equation}
    \hat{R}_z(\theta) = \hat{I} \cos\left(\frac{\theta}{2}\right) - i \sigma_z \sin\left(\frac{\theta}{2}\right)
\end{equation}
\begin{equation}
    \hat{R}_z(\theta) = \begin{pmatrix} 1 & 0 \\ 0 & 1 \end{pmatrix} \cos\left(\frac{\theta}{2}\right) - i \begin{pmatrix} 1 & 0 \\ 0 & -1 \end{pmatrix} \sin\left(\frac{\theta}{2}\right)
\end{equation}
\begin{equation}
    \hat{R}_z(\theta) = \begin{pmatrix} \cos\left(\frac{\theta}{2}\right) - i \sin\left(\frac{\theta}{2}\right) & 0 \\ 0 & \cos\left(\frac{\theta}{2}\right) + i \sin\left(\frac{\theta}{2}\right) \end{pmatrix}
\end{equation}
\begin{equation}
    \hat{R}_z(\theta) = \begin{pmatrix} e^{-i\theta/2} & 0 \\ 0 & e^{i\theta/2} \end{pmatrix}
\end{equation}

\begin{equation}
    \text{CNOT} = |0\rangle\langle 0| \otimes \hat{I} + |1\rangle\langle 1| \otimes \sigma_x
\end{equation}
\begin{equation}
    \text{CNOT} = \begin{pmatrix} 1 & 0 \\ 0 & 0 \end{pmatrix} \otimes \begin{pmatrix} 1 & 0 \\ 0 & 1 \end{pmatrix} + \begin{pmatrix} 0 & 0 \\ 0 & 1 \end{pmatrix} \otimes \begin{pmatrix} 0 & 1 \\ 1 & 0 \end{pmatrix}
\end{equation}
\begin{equation}
    \text{CNOT} = \begin{pmatrix} 1 & 0 & 0 & 0 \\ 0 & 1 & 0 & 0 \\ 0 & 0 & 0 & 0 \\ 0 & 0 & 0 & 0 \end{pmatrix} + \begin{pmatrix} 0 & 0 & 0 & 0 \\ 0 & 0 & 0 & 0 \\ 0 & 0 & 0 & 1 \\ 0 & 0 & 1 & 0 \end{pmatrix}
\end{equation}
\begin{equation}
    \text{CNOT} = \begin{pmatrix} 1 & 0 & 0 & 0 \\ 0 & 1 & 0 & 0 \\ 0 & 0 & 0 & 1 \\ 0 & 0 & 1 & 0 \end{pmatrix}
\end{equation}
sehingga CNOT dengan kontrol $q_0$ dan target $q_1$ adalah 
\begin{equation}
    \text{CNOT}_{q_0, q_1} = \begin{pmatrix} 1 & 0 & 0 & 0 \\ 0 & 1 & 0 & 0 \\ 0 & 0 & 0 & 1 \\ 0 & 0 & 1 & 0 \end{pmatrix}
\end{equation}
\begin{equation}
    \begin{split}
    |00\rangle &\to |00\rangle \\ 
    |01\rangle &\to |01\rangle \\ 
    |10\rangle &\to |11\rangle \\ 
    |11\rangle &\to |10\rangle
\end{split}
\end{equation}
namun CNOT dengan kontrol $q_1$ dan target $q_0$ adalah 
\begin{equation}
    \text{CNOT}_{q_1, q_0} = \begin{pmatrix} 1 & 0 & 0 & 0 \\ 0 & 0 & 0 & 1 \\ 0 & 0 & 1 & 0 \\ 0 & 1 & 0 & 0 \end{pmatrix}
\end{equation}
\begin{equation}
    \begin{split}
    |00\rangle &\to |00\rangle \\ 
    |01\rangle &\to |11\rangle \\ 
    |10\rangle &\to |10\rangle \\ 
    |11\rangle &\to |01\rangle
\end{split}
\end{equation}

\begin{equation}
    |\psi\rangle = \alpha |0\rangle + \beta |1\rangle
\end{equation}
\begin{itemize}
    \item $\alpha, \beta$ : Amplitudo probabilitas $(\alpha, \beta \in \mathbb{C})$.
\end{itemize}

\begin{equation}
    \langle \psi | \psi \rangle = 1
\end{equation}
\begin{equation}
    (\alpha^* \langle 0 | + \beta^* \langle 1 |) (\alpha |0\rangle + \beta |1\rangle) = 1
\end{equation}
\begin{equation}
    |\alpha|^2 \langle 0|0\rangle + \alpha^* \beta \langle 0|1\rangle + \beta^* \alpha \langle 1|0\rangle + |\beta|^2 \langle 1|1\rangle = 1
\end{equation}
\begin{equation}
    |\alpha|^2 + |\beta|^2 = 1
\end{equation}

\begin{equation}
    \hat{H} |\psi_n\rangle = E_n |\psi_n\rangle
\end{equation}
\begin{itemize}
    \item $\hat{H}$ : Operator Hamiltonian.
    \item $|\psi_n\rangle$ : Eigenstate ke-$n$.
    \item $E_n$ : Eigenvalue energi ke-$n$.
\end{itemize}

\begin{equation}
    E(\theta) = \frac{\langle \psi(\theta) | \hat{H} | \psi(\theta) \rangle}{\langle \psi(\theta) | \psi(\theta) \rangle}
\end{equation}
\begin{equation}
    E(\theta) = \langle \psi(\theta) | \hat{H} | \psi(\theta) \rangle
\end{equation}
\begin{itemize}
    \item $E(\theta)$ : Nilai ekspektasi energi (\textit{cost function}).
    \item $|\psi(\theta)\rangle$ : \textit{Trial ansatz} terparameterisasi.
\end{itemize}

\begin{equation}
    E(\theta) = \langle \psi(\theta) | \left( \sum_n E_n | \psi_n \rangle \langle \psi_n | \right) | \psi(\theta) \rangle
\end{equation}
\begin{equation}
    E(\theta) = \sum_n E_n |\langle \psi_n | \psi(\theta) \rangle|^2 \geq E_0
\end{equation}
\begin{itemize}
    \item Prinsip Variasional Rayleigh-Ritz: $E(\theta) \geq E_0$.
\end{itemize}

